% ==============================================================================
% reports/sections/results/06_07_xgboost_refit_study.tex
% ==============================================================================

\section{Nghiên cứu điển hình: Tác động của thiết kế pipeline và cờ refit trên XGBoost}
\label{sec:results_refit}

Trong quá trình đối soát thực nghiệm giữa các đường ống trong mã nguồn, nhóm phát hiện sự phân kỳ số liệu đáng chú ý đối với thuật toán XGBoost:
\begin{itemize}
    \item Trong benchmark chuẩn mực tại \path{artifacts/week3_week4/test_comparison.csv}, XGBoost đạt $\FOne = \textbf{0.0000}$ ($\text{SSH Recall} = 0.00\%$).
    \item Trong kết quả tuning độc lập tại \path{experiments/w3_05_time_tuning/test_eval/test_metrics.json}, XGBoost lại đạt $\FOne = \textbf{\XGBRefitTestFOne}$ ($\text{SSH Recall} = \textbf{\XGBRefitTestSSHRecall}$).
\end{itemize}

\subsection{Truy vết nguyên nhân kỹ thuật trong mã nguồn}
Tiến hành phân tích mã nguồn \path{src/ids/tune_xgboost.py}, nguyên nhân kỹ thuật được xác định tại dòng 151--168:
\begin{lstlisting}[language=Python, caption={Đoạn mã cấu hình RandomizedSearchCV và gộp dữ liệu gây kích hoạt refit trên tập gộp.}]
search = RandomizedSearchCV(
    estimator=base_model,
    param_distributions=param_distributions,
    n_iter=n_iter,
    scoring=f1_scorer,
    cv=cv_split, # Phân tách tùy biến giữa Train và Validation
    # refit=True là giá trị mặc định của scikit-learn
)
X_combined = pd.concat([X_train, X_val], axis=0).reset_index(drop=True)
y_combined = pd.concat([y_train, y_val], axis=0).reset_index(drop=True)
search.fit(X_combined, y_combined)
joblib.dump(search.best_estimator_, best_model_path)
\end{lstlisting}

Theo tài liệu đặc tả của thư viện \texttt{scikit-learn} \cite{scikit_learn_refit}, tham số \texttt{refit} có giá trị mặc định là \texttt{True}. Sau khi hoàn tất việc đánh giá các tổ hợp siêu tham số trên \texttt{cv\_split} (sử dụng tập con Train làm fold huấn luyện và Val làm fold kiểm định), đối tượng \texttt{RandomizedSearchCV} hoạt động đúng theo đặc tả: gọi hàm \texttt{fit()} của mô hình tốt nhất trên toàn bộ tập dữ liệu đầu vào đã truyền vào hàm \texttt{fit()}, tức là toàn bộ tập gộp $\mathbf{X}_{\text{combined}} = \Xtrain \cup \Xval$.

Lỗi ở đây nằm ở thiết kế đường ống: dữ liệu Train và Validation được chủ động gộp trước khi truyền vào \texttt{RandomizedSearchCV.fit()}, trong khi \texttt{refit=True} khiến mô hình tốt nhất được huấn luyện lại trên toàn bộ tập gộp đó.

\subsection{Hậu quả đối với bản chất bài toán đánh giá}
Trong khi tập $\Dtrain$ không chứa mẫu \SSH\ nào (0 flow), thì tập kiểm định $\Dval$ lại chứa \textbf{1,886 flow tấn công \SSH}.

Khi cơ chế \texttt{refit=True} diễn ra trên tập dữ liệu gộp:
\begin{enumerate}
    \item Mô hình XGBoost được huấn luyện trực tiếp trên 1,886 mẫu tấn công SSH từ tập Validation.
    \item Bài toán đánh giá trên tập Test theo thời gian đã bị biến đổi bản chất từ \textbf{kịch bản tổng quát hóa FTP $\to$ SSH (không có SSH trong tập huấn luyện)} thành \textbf{học có giám sát trên biến thể đã biết (\textit{Seen-Subtype Supervised Learning}) do được huấn luyện trên 1,886 mẫu SSH}.
    \item Việc được học trên 1,886 mẫu SSH này đã giúp XGBoost tiếp thu phân phối đặc trưng và cổng 22 của SSH, làm cho điểm số trên tập Test tăng từ \textbf{0.0000 lên \XGBRefitTestFOne}.
\end{enumerate}

Bảng \ref{tab:refit_comparison} tổng hợp sự đối lập giữa hai chế độ đánh giá này.

\begin{table}[htbp]
\centering
\caption{Đối chiếu hiệu năng của XGBoost giữa chế độ đánh giá tổng quát hóa chuẩn mực và chế độ tiếp xúc validation do thiết kế pipeline (Refit).}
\label{tab:refit_comparison}
\adjustbox{max width=\textwidth}{%
\begin{tabular}{lllrr}
\toprule
\textbf{Chế độ thực nghiệm} & \textbf{Dữ liệu thực tế fit mô hình} & \textbf{Mẫu SSH trong tập fit} & \textbf{Test F1-Score} & \textbf{Test SSH Recall} \\
\midrule
Chế độ 1: Tổng quát hóa chuẩn mực & Duy nhất Train ($\Dtrain$) & \DataTrainSSH\ mẫu & \textbf{\XGBTestFOne} & \textbf{\XGBTestSSHRecall} \\
Chế độ 2: Pipeline gộp tập (Refit) & Train + Val ($\Dtrain \cup \Dval$) & \DataValidationSSH\ mẫu & \textbf{\XGBRefitTestFOne} & \textbf{\XGBRefitTestSSHRecall} \\
\bottomrule
\end{tabular}%
}
\end{table}

Phát hiện này mang lại một bài học thực hành giá trị trong kỹ nghệ MLOps: việc chủ động gộp Train và Validation trước khi gọi hàm tìm kiếm siêu tham số kết hợp với hành vi mặc định \texttt{refit=True} của thư viện có thể làm thay đổi bản chất bài toán đánh giá nếu không được kiểm soát ranh giới dữ liệu một cách cẩn trọng.
